\scanpage{176}
\mspage{92}
% archive sheet number: 175
% lines 1-3
La démonstration se fait par un AQT fastidieux (j'avoue ne pas avoir poussé jusqu'au dernier détail \ldots).

% lines 4-5
Pour terminer la démonstration de 4.4.\unsure{9|8}, il reste à prouver un dernier (?)

% lines 6-11, replacements above the cancelled hypotheses
\uline{Lemme 4.4.\unsure{11|10}}. Soit $C$ une \inserted{petite} catégorie, stable par \deleted{$\unsure{\varinjlim}$ \unsure{de type} $\unsure{J}$ de cardinal $\leq\pi$} \inserted{sommes de cardinal $\leq\pi$}. Alors $\operatorname{Ind}'_{\pi}C$ est stable \deleted{par les \unsure{mêmes} \unsure{sommes}} \inserted{\unsure{sommes} de cardinal $\leq\pi$} (et, pour mémoire, le foncteur $\alpha:C\mathrel{\rule{1.5em}{0.4pt}}\operatorname{Ind}'_{\pi}C$ y commute).

% lines 12-21
Appliquant \deleted{le lemme} \inserted{le cor.} 4.4.\unsure{8|9} ou le lemme 4.4.7, on est ramené à prouver la cofinalité des (4.4.6.3), quand $J$ est discret, $+$ le fait que $\underline{\operatorname{Hom}}(J,C)_{/\varphi}$ est grande devant $\pi$, dans ce cas. Mais la donnée de $\varphi$ revient à celle d'une famille $(\varphi_j)_{j\in J}$ d'objets de $\operatorname{Ind}'_{\pi}(C)$, et les foncteurs (4.4.6.3) \unsure{s'explicitent|s'explicite} comme les projections
\[
\prod_{j\in J}(C_{/\varphi_j})\unsure{\longleftrightarrow|\longrightarrow} C_{/\varphi_j}
\]
du produit total sur un des facteurs. Donc on est ramené à ceci:
